\documentclass[conference]{IEEEtran}

% --- Grundkonfiguration ---
\usepackage[utf8]{inputenc}
\usepackage[T1]{fontenc}
\usepackage[ngerman]{babel}

% --- Wissenschaftliche Pakete ---
\usepackage{cite}
\usepackage{url}
\usepackage{booktabs}

\begin{document}

\title{Erste Version der Einleitung:\\
Trade-offs der Virtualisierung auf x86-Systemen:\\
Eine quantitative Analyse von Performance-Overheads\\
und Side-Channel-Risiken}

\author{%
  \IEEEauthorblockN{Maximilian Wagner}
  \IEEEauthorblockA{%
    \textit{Studiengang IT-Sicherheit} \\
    \textit{Fachbereich Informatik und Medien} \\
    \textit{Technische Hochschule Brandenburg} \\
    maximilian.wagner@th-brandenburg.de
  }
}

\maketitle

\begin{abstract}
Dieses Dokument enthält die erste KI-assistiert erstellte Version der
Einleitung der wissenschaftlichen Hausarbeit zum Thema
mikroarchitektonischer Ressourcen-Interferenzen in virtualisierten
Umgebungen. Im Anhang sind der verwendete System-Kontext sowie der
eingesetzte Prompt zur Nachvollziehbarkeit des Erstellungsprozesses
dokumentiert.
\end{abstract}

% =============================================================================
\section{Einleitung}
% =============================================================================

\paragraph{Thematische Einführung}
Virtualisierung bezeichnet die softwaregestützte Abstraktion physischer
Rechenhardware in eine oder mehrere logisch isolierte
Ausführungsumgebungen, sogenannte Gastsysteme. Der Hypervisor vermittelt
dabei zwischen den Gastsystemen und der physischen Hardware und
gewährleistet, dass Gastsysteme weder den Adressraum noch die
Ausführungskontrolle benachbarter Instanzen direkt adressieren können.
Diese logische Isolation ist das tragende Sicherheitsprinzip der
Multi-Tenancy-Architektur. Auf mikroarchitektonischer Ebene ist diese
Isolation jedoch unvollständig: Ressourcen wie der Last Level Cache
(LLC), die Hauptspeicherbandbreite und I/O-Subsysteme werden zwischen
allen auf denselben physischen Kernen eingeplanten Gastsystemen physisch
geteilt. Der Hypervisor abstrahiert diese Ressourcen nicht weg, sondern
macht deren gemeinsame Nutzung für das Gast-Betriebssystem transparent
--- und entzieht ihr damit gleichzeitig jede Kontrollmöglichkeit
\cite{koh2007analysis}.

\paragraph{Relevanz und Gefährdungspotenzial}
Die unkontrollierte Teilung mikroarchitektonischer Ressourcen begründet
eine systemische Schwachstellenklasse, die in modernen
Cloud-Infrastrukturen mit wirtschaftlich signifikantem Schadenpotenzial
verbunden ist. Koh et al.\ zeigen, dass Virtualisierungstechnologien
zwar funktionale Fehler- und Sicherheitsisolation bereitstellen, eine
strikte Performance-Isolation jedoch strukturell nicht leisten können
\cite{koh2007analysis}. Die daraus resultierende Ressourcen-Konkurrenz
manifestiert sich in zwei Angriffsformen: Erstens ermöglicht die
unkontrollierte Verdrängung von Cache-Lines durch einen aggressiven
Workload --- den sogenannten Noisy-Neighbor --- eine ungezielte
Degradation der Latenz und des Durchsatzes koexistierender Gastsysteme
und damit einen mikroarchitektonischen Denial-of-Service. Zweitens
erlauben deterministische Zugriffsmuster auf den Shared LLC über
Techniken wie PRIME+PROBE die Rekonstruktion von Laufzeitverhalten
über logische Isolationsgrenzen hinweg, was als Seitenkanal-Angriff
klassifiziert wird \cite{ge2018survey}. Das National Institute of
Standards and Technology (NIST) klassifiziert die unkontrollierte
Ressourcenerschöpfung auf Hypervisor-Ebene als kritisches
Plattform-Risiko und fordert hardwaregestützte Isolationsmechanismen
als Gegenmaßnahme \cite{nist2014hypervisor}.

\paragraph{Zielsetzung und Problemstellung}
Die referenzierte Literatur belegt die prinzipielle Ausnutzbarkeit
geteilter mikroarchitektonischer Ressourcen auf Basis älterer
CPU-Generationen. Für aktuelle Hardware-Architekturen besteht eine
empirische Forschungslücke: Weder die hybriden Kern-Topologien
(P-Cores/E-Cores) noch die hardwaregestützten Isolationsmechanismen
moderner Intel-Prozessoren --- insbesondere die Intel Resource Director
Technology (RDT) mit Cache Allocation Technology (CAT) --- wurden in
diesem Kontext systematisch untersucht
\cite{koh2007analysis,ge2018survey}. Diese Arbeit verfolgt das Ziel,
folgende Hypothese empirisch zu falsifizieren: \textit{Historische
Interferenz-Vektoren, die auf der gemeinsamen Nutzung des Last Level
Caches beruhen, führen auf einer Intel Raptor Lake Architektur unter
produktiven Hypervisor-Bedingungen (Proxmox VE) nicht mehr zu statistisch
quantifizierbaren Isolationsbrüchen.} Die Forschungsfrage lautet
entsprechend: Sind LLC-basierte Noisy-Neighbor-Interferenzen auf
aktueller Hardware mit modernen Mitigierungen empirisch nachweisbar?

\paragraph{Lösungsansatz}
Die Forschungsfrage wird durch einen empirischen Messansatz unter
deterministisch kontrollierten Bedingungen adressiert. Anstelle einer
rein theoretischen Analyse werden Interferenzeffekte direkt am
Zielsystem provoziert und quantifiziert. Um Messrauschen durch
nicht-deterministische Hardware-Verhalten auszuschließen, werden auf dem
Host-System tiefe C-States sowie dynamische Frequenzskalierung (Intel
Turbo Boost, CPU-Governor) deaktiviert. Die Messungen erfassen Durchsatz
und Latenz des Opfer-Systems sowohl unter isolierter Baseline-Last als
auch unter gleichzeitiger Störlast durch ein Angreifer-System. Das Delta
beider Messreihen quantifiziert den Isolationsbruch.

\paragraph{Praktische Demonstration}
Im Rahmen des Proof of Concept (PoC) werden zwei Debian-Gastsysteme
--- ein Angreifer und ein Opfer --- über striktes CPU-Pinning
(CPU Affinity) auf identische physische P-Cores des Proxmox-Hosts
fixiert, um die erzwungene LLC-Konutzung sicherzustellen. Das
Angreifer-System sättigt den LLC gezielt durch
\texttt{stress-ng {-}{-}cache 2 {-}{-}cache-level 3}. Das
Opfer-System erfasst parallel mittels \texttt{sysbench} und \texttt{fio}
den I/O-Durchsatz (MB/s) sowie die 95.~Perzentil-Latenz (ms) als
primäre Interferenzmetriken. Sollten aktuelle Hardware-Mitigierungen
--- insbesondere Intel RDT/CAT --- die Interferenz auf ein statistisch
nicht signifikantes Maß reduzieren und die Hypothese damit nicht
falsifizierbar machen, wird eine methodische Fallback-Strategie
aktiviert: Ein systematischer Leistungsvergleich von reiner Emulation
(QEMU), Para-Virtualisierung (LXC) und hardwareunterstützter
Virtualisierung (KVM) unter identischer Last und strikter
Host-Kontrolle. Der Fallback liefert in diesem Fall unabhängige,
empirisch verwertbare Aussagen über den inhärenten Performance-Overhead
der jeweiligen Virtualisierungsparadigmen.

\paragraph{Aufbau der Arbeit}
Abschnitt~II systematisiert den Stand der Forschung zu
mikroarchitektonischen Interferenz-Vektoren, Virtualisierungsparadigmen
und hardware-seitigen Mitigierungsstrategien. Abschnitt~III beschreibt
den Versuchsaufbau, die Host-Konfiguration sowie die eingesetzten
Messwerkzeuge und -methoden. Abschnitt~IV präsentiert und bewertet die
Ergebnisse des PoC respektive des Fallback-Benchmarks. Abschnitt~V
diskutiert die Befunde im Kontext der Ausgangshypothese und der
referierten Literatur. Abschnitt~VI fasst die zentralen Erkenntnisse
zusammen und benennt weiterführende Forschungsfragen.

\bibliographystyle{IEEEtran}
\bibliography{../references}

% =============================================================================
% Ab hier: Anhang in Einzelspalte
% =============================================================================
\clearpage
\onecolumn

\section*{Anhang: Erstellungs-Dokumentation}

\noindent
Dieser Anhang dokumentiert den KI-gestützten Erstellungsprozess gemäß
den Anforderungen der Lehrveranstaltung. Das verwendete KI-Werkzeug
ist Claude Sonnet~4.6 (Anthropic, Stand: Mai 2026).

% --- Kontext -------------------------------------------------------------------
\subsection*{A. Verwendeter System-Kontext}

\noindent
Der folgende System-Kontext wurde dem KI-Werkzeug als persistente
Instruktionsgrundlage übergeben. Er legt Projektziel, Hardware-Szenario,
Methodik und Fallback-Strategie fest.

\bigskip
\noindent\textbf{A.1\quad Projektziel und Rahmenbedingungen}

\begin{itemize}
  \item \textbf{Format:} IEEE Conference Format
        (\texttt{\textbackslash documentclass[conference]\{IEEEtran\}})
  \item \textbf{Sprache:} Deutsch (wissenschaftlicher Standard)
  \item \textbf{Kernaufgabe:} Untersuchung der Isolation in
        Virtualisierungslösungen mit Fokus auf Hardware-Sicherheit
\end{itemize}

\bigskip
\noindent\textbf{A.2\quad Erweiterte Aufgabenstellung (Fokus: Noisy-Neighbor)}

Das Projekt erweitert die ursprüngliche Basisaufgabe (Vergleich von
Virtualisierungsparadigmen) um eine spezifische Untersuchung
mikroarchitektonischer Ressourcen-Interferenzen. Ziel ist die Falsifikation
der Hypothese, dass historische Angriffsvektoren auf moderner Hardware
(Intel Raptor Lake) durch aktuelle Isolationsmechanismen
(Intel RDT/CAT) vollständig neutralisiert werden.

\bigskip
\noindent\textbf{A.3\quad Hardware-Szenario und Host-Konfiguration}

\begin{itemize}
  \item \textbf{Host:} Proxmox VE auf Intel Raptor Lake
        (hybride P-Core/E-Core Architektur)
  \item \textbf{Shared Resource:} Last Level Cache (LLC),
        physisch geteilt von P-Cores
  \item \textbf{Host-Determinismus (zwingend):}
    \begin{itemize}
      \item Deep C-States deaktivieren (Kernel-Parameter):\\
            \texttt{intel\_idle.max\_cstate=1
            processor.max\_cstate=1 idle=poll}
      \item Intel Turbo Boost deaktivieren:\\
            \texttt{echo 1 >\
            /sys/devices/system/cpu/intel\_pstate/no\_turbo}
      \item CPU-Governor auf \texttt{performance} fixieren
    \end{itemize}
\end{itemize}

\bigskip
\noindent\textbf{A.4\quad Methodik und Proof of Concept (PoC)}

\begin{itemize}
  \item \textbf{Setup:} Zwei identische Debian-Instanzen
        (Attacker \& Victim)
  \item \textbf{Isolations-Bruch:} CPU-Pinning beider Instanzen auf
        denselben physischen P-Core (CPU Affinity in Proxmox)
  \item \textbf{Metriken:} Durchsatz (MB/s) und 95.~Perzentil-Latenz (ms)
        via \texttt{sysbench} und \texttt{fio} auf dem Victim-System;
        Attacker läuft \texttt{stress-ng {-}{-}cache 2
        {-}{-}cache-level 3}
\end{itemize}

\bigskip
\noindent\textbf{A.5\quad Fallback-Strategie}

Greift, wenn keine messbaren Interferenzen auftreten:

\begin{itemize}
  \item \textbf{Thema:} Vergleich von Emulation (QEMU),
        Para-Virtualisierung (LXC) und hardwareunterstützter
        Virtualisierung (KVM)
  \item \textbf{Fokus:} Leistungsfähigkeit bei Rechen- und
        I/O-Operationen unter strikter Host-Kontrolle
\end{itemize}

% --- Prompt -------------------------------------------------------------------
\bigskip
\subsection*{B. Verwendeter Prompt}

\noindent
Der folgende Prompt wurde zur Generierung der Einleitung eingesetzt.
Das KI-Werkzeug erhielt diesen Prompt zusammen mit dem unter A.
dokumentierten System-Kontext sowie den Quelldateien
\texttt{expose\_hardware\_security.tex},
\texttt{literatur\_recherche\_draft.tex} und \texttt{references.bib}.

\bigskip
\begin{quote}
\noindent\textbf{Rollenprofil:} Wissenschaftlicher Autor im Fachbereich
IT-Sicherheit und Hardware-Architektur.

\medskip
\noindent\textbf{Aufgabe:} Erstellung eines inhaltlichen Entwurfs der
Einleitung (Ausgabe als Markdown-Dokument \texttt{Einleitung.md}) für
eine wissenschaftliche Hausarbeit im IEEE-Format. Die inhaltliche
Struktur muss strikt und sequenziell den folgenden sechs Leitfragen
folgen.

\medskip
\noindent\textbf{Eingabedaten (zwingend zu berücksichtigen):}
\begin{itemize}
  \item System-Kontext: \texttt{ki\_agent\_context\_projekt.md}
  \item Projektskizze:
        \texttt{project/expose/expose\_hardware\_security.tex}
  \item Stand der Forschung:
        \texttt{paper/literatur\_recherche\_draft.tex}
  \item Literaturdatenbank: \texttt{paper/references.bib}
\end{itemize}

\medskip
\noindent\textbf{Inhaltliche Strukturierung (Absatzweise abzuarbeiten):}
\begin{enumerate}
  \item \textbf{Worum geht es? (Thematische Einführung):} Sachliche
        Einführung in Virtualisierung, logische Isolation und die
        physische Teilung mikroarchitektonischer Ressourcen
        (insbesondere Caches).
  \item \textbf{Warum ist das Thema relevant? (Gefährdungspotenzial \&
        Aktualität):} Einordnung der Relevanz für moderne
        Cloud-Infrastrukturen. Der Noisy-Neighbor-Effekt als
        systemische Schwachstelle (Leistungsdegradation,
        Timing-Angriffe, Denial-of-Service auf Hardware-Ebene).
  \item \textbf{Was ist das Ziel dieser Hausarbeit? (Zielsetzung \&
        Problemstellung):} Falsifikation der Hypothese, dass
        historische Interferenz-Vektoren auf aktuellen,
        hardware-mitigierten Architekturen (Intel Raptor Lake) unter
        produktiven Hypervisor-Bedingungen (Proxmox VE) weiterhin
        quantifizierbare Isolationsbrüche verursachen.
  \item \textbf{Was ist der Lösungsansatz?:} Empirische
        Quantifizierung der Performance-Isolation unter kontrollierter
        Störlast.
  \item \textbf{Was wird im praktischen Teil demonstriert?:}
        Proof of Concept: CPU-Pinning auf physische P-Cores,
        LLC-Sättigung via \texttt{stress-ng}, Baseline- und
        Degradationsmessung via \texttt{sysbench} und \texttt{fio}.
        Zwingende Integration der Fallback-Strategie (Leistungsvergleich
        QEMU\,/\,LXC\,/\,KVM), falls Intel RDT den Isolationsbruch
        neutralisiert.
  \item \textbf{Wie ist die Gliederung der Arbeit?:} Rein deskriptiver
        Ausblick auf die nachfolgenden Kapitel (Stand der Forschung,
        Methodik, Evaluation, Fazit).
\end{enumerate}

\medskip
\noindent\textbf{Stilistische und formale Vorgaben:}
\begin{itemize}
  \item Sprache: Wissenschaftliches, objektives Deutsch
  \item Terminologie: Korrekte Fachsprache (IT-Security\,/\,Hardware-%
        Architektur), keine unnötigen Vereinfachungen
  \item Wissenschaftlicher Standard: Differenzierung zwischen Fakt
        (Architekturgegebenheiten) und Hypothese
  \item Formatierung: Ausgabe als validierbarer Markdown-Codeblock,
        keine Meta-Kommentare
  \item Referenzierung: Platzhalter im \texttt{\textbackslash
        cite\{\}}-Format, korrespondierend zu den Einträgen
        in der \texttt{references.bib}
\end{itemize}
\end{quote}

\end{document}
